\section{One binomial certificate for factorial and CRT arithmetic}
\label{sec:direct319}

The strongest direct result of this investigation is a literal quartic in
319 nonnegative integers. It improves the latest inspected direct natural
equation in four measured quantities at once:
\[
\begin{array}{lrrrr}
 &\text{variables}&\text{quadratic rows}&\text{monomials}&\text{height}\\
\text{Round 14 derived-label source}&349&268&12\,768&86\,915\,503\\
\text{Merged factorial--CRT source}&319&254&10\,283&86\,842\,335.
\end{array}
\]
Both single equations have exact degree four. The comparison concerns
direct C222 equations with completely collected coefficients. It does not
assert superiority to the repository's few-variable equations with different
coefficient and degree costs.

\begin{theorem}[Explicit direct equation]\label{thm:direct319}
Let $R_1,\ldots,R_{254}\in\Z[z_0,\ldots,z_{318}]$ be the complete quadratic
residual list in the delivered file
\texttt{direct/equation319/natural/equation.json}, and put
\[
             P_{319}(z)=\sum_{j=1}^{254}R_j(z)^2.
\]
The supplied collected coefficient file defines the same polynomial.
Its degree is four, it has $10\,283$ nonzero monomials, and its height is
$86\,842\,335$. With the inherited numerical-calculus calibration,
\[
 \mathrm{ZFC}\vdash
 \bigl(\nexists z\in\N^{319}:P_{319}(z)=0\bigr)
       \ \longleftrightarrow\ \operatorname{Con}(\mathrm{ZFC}).
\]
Here $\N=\{0,1,2,\ldots\}$; all coefficients are fixed, literal integers.
\end{theorem}

The ordinary proof is given below. Its arithmetic source has 465 natural
coordinates and 400 quadratic rows; an independent exact certificate
eliminates 146 coordinates, including a nonnegativity proof for every
restoration. Expansion of all remaining squares verifies every final
coefficient. These finite checks support the literal construction; the
unbounded equivalence is proved mathematically, rather than inferred from
tests. No new end-to-end Lean proof term or emitted object-language ZFC
derivation is claimed.

\subsection{The improvement and its mathematical provenance}

The preceding reciprocal-CRT construction used $F=t!$ and
\[
 B=\binom{F^2+t}{t},
\]
but obtained the factorial through a different binomial,
$\binom{Ft^2+t}{t}$. We align these two uses of binomial arithmetic.
The squared-radius factorial fixed-point theorem was already established
in package 105, \emph{Bootstrap certificates}, Section 2 of its
\texttt{NEW\_RESULTS.md}; it is not a new theorem of this iteration.
The common Pell-residue representation was already established in package
108, \emph{Pell residues instead of power values}. The shared-factorial
source, reciprocal prime-power CRT and derived-label range come from the
latest Round 14 packet. These source proofs are retained in the archive.

The completed contribution is their integration: the \emph{same} binomial
output serves the factorial and CRT product; an independent radix bound
validates that output before either semantic interpretation is assumed;
common Pell residues eliminate eleven full numerator powers; and the
entire construction is compiled, contracted and measured. This is a new
complete equation relative to the inspected snapshot, without a claim of
worldwide priority for the ingredients.

\subsection{The inherited factorial identity, with an independent proof}

\begin{lemma}[Squared-radius factorial]\label{lem:direct-factorial}
For $t\ge6$ and positive integer $F$,
\[
 F=\left\lceil\frac{(F^2+1)^t}{\binom{F^2+t}{t}}\right\rceil
       \quad\Longleftrightarrow\quad F=t!.
\]
\end{lemma}
\begin{proof}
Put
\[
 H_t(F)=t!\prod_{i=1}^{t}\frac{F^2+1}{F^2+i}.
\]
The ceiling assumption gives $H_t(F)\le F\le t!$. If $F<t$, monotonicity
of each factor in its radius gives
\[
 H_t(F)\ge\frac{2^t}{t+1}>t>F,
\]
a contradiction. The middle inequality follows at $t=6$ and persists
by induction. Thus $F\ge t$.

Using $\prod(1-u_i)\ge1-\sum u_i$,
\[
 \prod_{i=1}^{t}\frac{F^2+1}{F^2+i}
 >1-\frac{t(t-1)}{2F^2}>\frac12.
\]
If $F<t^2$, then $H_t(F)>t!/2>t^2>F$, again impossible. Here
$t!>2t^2$ for $t\ge6$ follows immediately by induction. Hence $F\ge t^2$,
and the stronger estimate is
\[
 H_t(F)>t!\left(1-\frac1{2F}\right).
\]
If $t!\ge F+1$, the right side is at least
$(F+1)(1-1/(2F))=F+1/2-1/(2F)\ge F$, contradicting $H_t(F)\le F$.
Therefore $F=t!$. Conversely, for $F=t!$, the same estimate gives
$F-1/2<H_t(F)\le F$, so its ceiling is $F$.
\end{proof}

Package 105 proves the more general statement for every nonnegative input.
The restricted proof above is enough here because the horizon is padded
to at least 16. The ceiling has an exact natural graph:
\begin{equation}\label{eq:direct-ceiling}
 V+B=BF+r+1,\qquad B=r+s+1,\qquad r,s\in\N.
\end{equation}
For $B>0$ these rows give $V+s=BF$ and $0\le s<B$, hence
$F=\lceil V/B\rceil$. If $V\ge1$ they also force $F\ge1$.

\subsection{A radix that can be justified before the factorial}

Define, using natural variables $t_0,x,y,F,b_0$,
\begin{equation}\label{eq:direct-radix}
 \begin{gathered}
 t=t_0+16,\qquad R=2+t+x+y+F,\qquad K=2R,\\
 B=b_0+2,\qquad N=B^K,\qquad q=N^K,\qquad V=(F^2+1)^t.
 \end{gathered}
\end{equation}
These variable powers are explanatory requests: their final implementation
uses the explicit indexed-Pell graph below. Once they have been recovered,
the coordinate shift $B\ge2$ gives
\begin{equation}\label{eq:direct-independent-radix}
 N\ge2^K,\qquad q=B^{K^2}\ge2^{4R^2}.
\end{equation}
No interpretation of $B$ as a binomial is needed for this bound.

For $n=F^2+t$ and $0\le j\le t$,
\begin{equation}\label{eq:direct-first-digit-bound}
 \binom nj\le n^t\le(R^2+R)^t
       \le2^{2Rt}<2^{2R^2}<q.
\end{equation}
Use the zero convention if $j>n$. The elementary bound $R^2+R\le4^R$
holds for $R\ge1$, and $t<R$. Thus the first exact digit graph can be
validated \emph{before} any factorial or CRT meaning is available. It yields
$B=\binom{F^2+t}{t}$; equations~\eqref{eq:direct-ceiling} and
Lemma~\ref{lem:direct-factorial} then yield $F=t!$.

This order removes the potential circularity in identifying the two
binomials. The previous radix $N(N+V+1)^K$ is replaced by $N^K$; its
bootstrap is supplied by an explicit lower bound on $B$, not by assuming
the value of the very digit being tested.

\subsection{Six exact powers and eleven modular Pell residues}

For $a\ge2$, write
\[
 X_0(a)=1,\quad Y_0(a)=0,\quad
 X_{j+1}=aX_j+(a^2-1)Y_j,\quad Y_{j+1}=X_j+aY_j.
\]
The retained compiler supplies indexed-Pell certificates for these pairs,
sharing an auxiliary Pell solution and CRT base across all requested
indices. The final system uses exactly six full power requests:
\begin{equation}\label{eq:direct-six-powers}
 \begin{aligned}
 V&=(F^2+1)^t,& N&=B^K,&q&=N^K,\\
 D_t&=q^t,&D_y&=q^{y+1},&D_x&=q^{x+1}.
 \end{aligned}
\end{equation}
The denominator $D_{t+1}=qD_t$ is obtained by multiplication, so it needs
no separate power request or Pell pair at index $t+1$.

Let $W$ be two plus the sum of the distinct bases, exponents and outputs
in~\eqref{eq:direct-six-powers}. The common frame is
\begin{equation}\label{eq:direct-free-frame}
 L^2-((W+1)^2-1)(WZ)^2=1,\qquad Z\ge1,\qquad A=L+\sigma.
\end{equation}
All coordinates are natural. Pell classification and
$Y_e(W+1)\equiv e\pmod W$ imply $L>W^W$, hence $A>W^W$. Conversely the
index $e=W$ supplies $L,Z$, after which every $A\ge L$ is allowed.

For an exact request $m=n^k$, the compiler asserts
\[
 T_n=2An-n^2-1\ge0,\qquad
 X_k(A)+nY_k(A)=AY_k(A)+m+T_nh,\quad h\ge0.
\]
The first row excludes $n=0$. The frame bounds $m$ and $n^k$ strictly
below $T_n$, and the Pell recurrence gives their congruence modulo $T_n$,
so $m=n^k$. Every exponent is positive by its literal definition.
This is the inherited simultaneous positive-power theorem; separating
the lower bound $L$ from the actual base $A$ permits the congruence below.

Put
\begin{equation}\label{eq:direct-root-modulus}
 b=q+1,\qquad \mathcal M=qD_{t+1}D_xD_y
                           =q^{t+x+y+4},
\end{equation}
and impose one natural-quotient equation
\begin{equation}\label{eq:direct-common-root}
              2A(q+1)=(q+1)^2+1+\mathcal M h_0.
\end{equation}
For every $n\ge0$ it yields
\begin{equation}\label{eq:direct-residue}
 \mathcal R_n=X_n(A)+(q+1-A)Y_n(A)
                    \equiv(q+1)^n\pmod{\mathcal M}.
\end{equation}
Indeed the substitution $s=q+1-A$ respects $s^2=A^2-1$ modulo
$\mathcal M$, and $(A+s)^n=X_n+Y_ns$. Moreover $\mathcal R_n>0$:
for $n>0$, $X_n>(A-1)Y_n$, and for $n=0$ the value is 1.

For every required denominator $d$, both $qd$ and $Nd$ divide
$\mathcal M$: the first follows from its product definition and the
second from $N\mid q$. The one exact binomial uses
\begin{equation}\label{eq:direct-exact-digit}
 \begin{aligned}
 X_n+(q+1)Y_n&=AY_n+r+dB+qd h,\\
 d&=r+s+1,&q&=B+o+1,
 \end{aligned}
\end{equation}
where $n=F^2+t$, $d=D_t$. Every divisibility request
$N\mid\binom nk$ instead uses
\begin{equation}\label{eq:direct-div-digit}
 X_n+(q+1)Y_n=AY_n+r+Nd h,\qquad d=r+s+1,
 \qquad d=q^k.
\end{equation}
The binomial theorem identifies the digit or its divisibility once all
lower coefficients are below $q$. Positivity of $\mathcal R_n$ ensures
that the constructed quotients are nonnegative.

Completeness of the common root is not automatic and must be checked.
As $\mathcal M$ is a power of $q$,
$\gcd(q+1,\mathcal M)=1$. Thus $\gcd(2(q+1),\mathcal M)$ is 1 or 2;
in the second case $(q+1)^2+1$ is even. Equation~\eqref{eq:direct-common-root}
therefore defines a soluble congruence class for $A$, containing arbitrarily
large positive representatives. Choose one above $L$ and large enough to
make $h_0\ge0$. Then supply the standard common-index witnesses at this
base. All semantic values were fixed beforehand.

There are eleven modular requests: one exact binomial and ten ranges.
Their upper indices are positive and need not be included in $W$;
the indexed-Pell graph certifies those indices directly. Fifteen independent
indices suffice for the six exact powers and eleven residues together.

\subsection{The explicit numerical source and its finite proof flow}

For completeness, the arithmetic source is specified here before applying
the inherited calibration. Its constructors are
\[
 \begin{aligned}
 I(a,b)&=2(a+b)(a+b+1)+4b,& M(a,b)&=I(a,b)+1,\\
 A(a,b)&=I(a,b)+2,&E(a,b)&=I(a,b)+3,\\
 N_0(a)&=I(a,1),&X_0(a,b)&=N_0(A(a,N_0(b))).
 \end{aligned}
\]
The subscript here distinguishes the formula constructors from the global
CRT modulus $N$ and Pell abscissas. Define
\[
 \begin{aligned}
 \operatorname{req}(a,k,m)
   &=A(a,I(m,N_0(A(a,I(m,N_0(M(a,k))))))),\\
 s(a,b,k)&=A(k,I(M(k,a),M(k,b))),\\
 u(a,b,k,m)&=N_0(A(a,I(s(a,b,k),\operatorname{req}(b,k,m)))),\\
 \operatorname{clos}(a,b,k,m)
   &=A(b,I(M(b,k),N_0(A(a,I(m,u(a,b,k,m)))))),\\
 v(a,b,c)&=I(M(a,b),A(c,M(a,b))),\\
 h(a,b,c)&=A(c,I(\operatorname{clos}(a,b,c,c),
                          N_0(\operatorname{req}(b,c,c)))),\\
 \operatorname{bundle}(a,b,c)&=N_0(I(v(a,b,c),h(a,b,c))).
 \end{aligned}
\]
The nine conclusion polynomials are
\begin{align*}
 c_1&=X_0(a,E(a,b)),\\
 c_2&=I(A(a,I(b,c)),I(A(a,b),A(a,c))),\\
 c_3&=I(E(a,b),I(M(c,a),M(c,b))),\\
 c_4&=I(A(a,A(b,c)),A(b,A(a,c))),\\
 c_5&=I(E(c,a),I(M(c,b),M(a,b))),\\
 c_6&=I(X_0(a,A(a,b)),b),\\
 c_7&=I(E(c,a),I(E(c,b),A(M(a,b),E(a,b)))),\\
 c_8&=I(I(I(c,d),a),I(I(a,c),I(b,c))),\\
 c_9&=\operatorname{bundle}(a,b,c).
\end{align*}
Their code definitions are matched to the retained C222 calculus.
The first eight are unconditional axiom families. The ninth is admitted
when $c\ne a,b$, represented by
$G=(c-a)^2(c-b)^2-(d+1)^2=0$.

Write $w(z)=1+Fz$. For the eight local columns
$(a,b,c,d,S,S',u,v)$ define
\begin{align*}
 H={}&(S'-S)^2\prod_{j=1}^{8}(S'w(c_j)-S)^2\\
 &\cdot\bigl((S'w(c_9)-S)^2+G^2\bigr)\\
 &\cdot\bigl(S'w(b)-Sw(a)w(I(a,b))\bigr)^2
       \bigl(S'w(A(a,b))-Sw(b)\bigr)^2,\\
 \mathcal L(i)&=H+\bigl(x-S-(1+(i+1)y)u\bigr)^2
                  +\bigl(x-S'-(1+(i+2)y)v\bigr)^2,\\
 J&=\sum_{j=1}^{9}c_j+a+b+I(a,b)+A(a,b)+1.
\end{align*}
The local bounds are $(t,t,t,t,y,y,x,x)$, together with the derived
condition $J\le t$. The endpoint equations are
\[
 x=F+1+(y+1)u_0,\qquad y=F+1+r_0,\qquad
 x=1+(1+(t+1)y)u_t.
\]
Every factor of $H$ is nonnegative. A zero selects truth, a genuine
axiom, modus ponens, or generalization, and gives the appropriate
accumulator identity
\[
 S'w(\text{conclusion})=Sw(\text{premise}_1)w(\text{premise}_2).
\]
Unused premises have truth label 0 and weight one.

Because $F=t!$, the weights $1+jF$ for $1\le j\le t$ are pairwise
coprime: a prime common to weights $j\ne k$ would divide $j-k$, hence
divide $F$, which is impossible. The beta equations identify overlapping
accumulators by uniqueness of remainders bounded by $y$. Their positive
initial value and positive weights propagate positivity. Multiplying the
local identities and cancelling gives
\[
 \prod w(\text{conclusions})=w(1)\prod w(\text{premises}).
\]
Taking valuations at a prime of each positive label's weight gives a
multiplicity excess of one at label 1 and zero at every other positive
label. If 1 were outside the derivable closure, each selected rule with
conclusion outside it would have a premise outside it. Counting these
occurrences contradicts that excess. This argument allows cyclic or
repeated selected rules; it does not assume the supplied flow is a tree.

Conversely expand a finite refutation to a proof tree, choose the horizon
above its length, parameters and all finite $J$-values, and pad with truth
steps. Define $F=t!$ only after this choice. Expanding the pending proof
multiset backwards gives the accumulator products. Choose $y$ large and
divisible by $t!$, then choose $x$ by CRT for the beta moduli. This is the
inherited complete proof-flow construction, with no dependence of the
horizon on the later accumulator sizes.

\subsection{Reciprocal prime-power decoding and the last carry bounds}

With $F=t!$, let
\[
 d_i=1+F^2/i\quad(1\le i\le t),\qquad
 B=\prod_{i=1}^{t}d_i=\binom{F^2+t}{t},\qquad N=B^K.
\]
The $d_i$ are integral, pairwise coprime, and have no prime divisors at
most $t$. These facts follow from $i\mid F$ and
$id_i-jd_j=i-j$. Introduce nine positive encoded representatives,
the eight local columns $A_j$ and index $C$, satisfying
\begin{equation}\label{eq:direct-crt-system}
 \begin{gathered}
 C+\sum_j A_j+g=9N^2,\qquad B\mid C+F^2+1,\\
 N\mid\binom Ct,\qquad N\mid\binom{A_j}{b_j+1},\qquad
 N\mid\mathcal L(C,A),\qquad N\mid\binom{J(A)}{t+1}.
 \end{gathered}
\end{equation}
The $b_j$ are the local bounds. The derived value $J(A)$ is deliberately
excluded from the aggregate bound.

The elementary one-sided range lemma states
\[
             p^e\mid\binom ak\quad\Longrightarrow\quad
                         a\bmod p^e<k.
\]
For $a\ge k$, the $p$-valuation of the binomial is the sum of
$\lfloor a/p^j\rfloor-\lfloor(a-k)/p^j\rfloor-\lfloor k/p^j\rfloor$,
each zero or one. If the interval $[a-k+1,a]$ had no multiple of $p^e$,
all terms from $j=e$ onward would vanish, contradicting valuation at
least $e$. The case $a<k$ is immediate. Its converse is false and is not
used.

Choose a prime $p\mid d_i$, write $m=v_p(d_i)$ and $e=Km$. The index
congruence and $p>t$ identify $C\bmod p^e=i-1$. The other ranges give
bounded residues for the local columns. A retained literal majorant
certificate gives $\mathcal L,J<R^E$, where $E=2^{23}$. Since
$R\ge16!>2^{44}$, elementary induction gives
$R^E<2^R<2^K\le p^e$. Polynomial congruence now yields
$\mathcal L(i-1,\text{local})=0$ and $J(\text{local})\le t$.

For completeness choose the representatives by CRT at precision
$d_i^{2K}$, so they lie in $[1,B^{2K}]=[1,N^2]$. Their falling factorial
numerators have a factor of valuation at least $2Km$; division by $k!$
loses at most $k-1\le K-1$. This proves the required divisibilities by
$p^{Km}$, including the derived polynomial range. Precision $2K$ is
essential to this argument; no false converse to the range lemma is used.

The remaining digit bounds also remain strict for the new radix $q=N^K$.
For $n\le9N^2$, $1\le k\le R-1$, and $j\le k$,
\[
 \binom nj\le n^k\le9^kN^{2k}
      <2^{4R}N^{2k}\le N^{2k+2}\le N^K=q.
\]
For the derived argument the positive-DAG certificate gives
$J(a,b,c,d)\le(n+1)^\ell$ when the four inputs are at most $n$,
with $\ell=2^{22}$. Since $n\le9N^2$ and $N\ge10$, $J(A)\le N^{3\ell}$.
The finite inequality $2\cdot16!\ge3\ell\cdot17$ persists by induction,
so $3\ell(t+1)\le2F<K$. Therefore
\[
 \binom{J(A)}j\le J(A)^{t+1}
                  \le N^{3\ell(t+1)}<N^K=q.
\]
This last strict inequality is necessary: using the old statement
$N^K<q$ would be incorrect for the new radix.

\subsection{Closing the theorem and reproducing the finite object}

Soundness first recovers the six exact powers from the independent Pell
frame. It then uses~\eqref{eq:direct-first-digit-bound} to interpret the
single exact binomial and the factorial ceiling, before invoking any CRT
or proof-label meaning. Only afterward are the ten ranges interpreted,
the bounded local source decoded, and the Horn-flow theorem applied.
The inherited ordinary calibration is
\[
 \mathrm{C222}(1)\ \longleftrightarrow\
 \text{finite ZF refutation}\ \longleftrightarrow\
 \text{finite ZFC refutation}.
\]
Its named interfaces are
\begin{quote}\small\ttfamily
KEndpoint.c222\_iff\_zf\_refutation\\
ConstructibleChoice.zf\_refutation\_iff\_zfc\_refutation.
\end{quote}
The construction uses the numerical-calculus endpoint before the
machine-specific theorem; it does not simulate K215 or instantiate a
Jones universal polynomial.

Completeness chooses, in order, the finite proof and padded horizon;
the factorial, proof-flow and beta data; $B,R,K,N,q$ and the CRT
representatives; the six powers and eleven true binomial residues; and
finally a lower Pell frame and sufficiently large $A$ in the common-root
congruence class. The common-index theorem supplies all auxiliary Pell
witnesses. Every displayed natural quotient is then reconstructed.
The checked natural projection gives a zero of $P_{319}$. Conversely
its restoration certificate returns every deleted coordinate, so it
preserves both existence directions.

The optional square-shell companion has 320 natural variables, 255
quadratic rows, $13\,828$ terms and coefficient height $15\,248\,812$.
It trades one extra variable and more terms for smaller coefficients,
using the retained constructor-preserving syntax transport. The main
319-variable result stays with the original Cantor encoding.

The complete source bundle uses only the Python 3 standard library.
Normal and optimized-interpreter replays rebuild both equations and
compare all twenty generated JSON files byte for byte. The independent
checker verifies every restoration, every final coefficient, all live
coordinates and deliberately corrupted controls. Separate finite
regressions cover $18\,049$ factorial cases, 270 independent radix
bounds, 49 growth budgets, $1\,530$ Pell-residue instances and 516 digit
and divisibility instances. Those checks are arithmetic regressions,
not proofs of unbounded statements or searches for a ZFC refutation.
